\documentclass[11pt,a4paper]{article}

\usepackage[utf8]{inputenc}
\usepackage[T1]{fontenc}
\usepackage[spanish,es-tabla]{babel}
\usepackage{mathptmx}
\usepackage[margin=2.5cm]{geometry}
\usepackage{booktabs}
\usepackage{tabularx}
\usepackage{xcolor}
\usepackage{fancyhdr}
\usepackage{hyperref}

\definecolor{azulOscuro}{RGB}{30, 58, 95}
\definecolor{azulClaro}{RGB}{2, 132, 199}
\definecolor{grisFondo}{RGB}{240, 244, 248}

\pagestyle{fancy}
\fancyhf{}
\fancyhead[L]{\small\textsf{Programación Web y Sistemas Inmersivos | UPV}}
\fancyhead[R]{\small\textsf{Documento de Relación de Trabajo}}
\fancyfoot[C]{\thepage}
\renewcommand{\headrulewidth}{0.5pt}
\renewcommand{\footrulewidth}{0.5pt}

\hypersetup{
    colorlinks=true,
    linkcolor=azulOscuro,
    urlcolor=azulClaro,
    pdftitle={Documento de Relacion de Trabajo - Tu Nevera Sostenible}
}

\begin{document}

% ==============================================================================
% PÁGINA 1: ALEJANDRO
% ==============================================================================
\begin{center}
    {\Large \textbf{UNIVERSITAT POLITÈCNICA DE VALÈNCIA}}\\[0.2cm]
    {\large Programación Web y Sistemas Inmersivos --- Curso 2026--2027}\\[0.3cm]
    {\LARGE \textbf{\color{azulOscuro}Documento de Relación de Trabajo Individual}}\\[0.2cm]
    {\small Proyecto: \textit{Tu Nevera Sostenible (ODS 2 y ODS 12)}}
\end{center}

\vspace{0.3cm}
\hrule
\vspace{0.4cm}

\subsection*{1. Datos del Estudiante}
\begin{tabularx}{\textwidth}{l X l X}
    \textbf{Nombre:} & Alejandro [Apellidos] & \textbf{DNI/NIE:} & [Completar DNI] \\
    \textbf{Email UPV:} & [alejandro@alumni.upv.es] & \textbf{Grupo:} & [Completar grupo] \\
\end{tabularx}

\vspace{0.4cm}
\subsection*{2. Ficheros Bajo Responsabilidad Única}
En cumplimiento de la normativa docente, cada fichero de código cuenta con un responsable exclusivo.
\begin{table}[h!]
\centering
\small
\begin{tabularx}{\textwidth}{p{4cm} p{3.5cm} X}
\toprule
\textbf{Fichero de Código} & \textbf{Tecnología} & \textbf{Descripción de Responsabilidad} \\
\midrule
\texttt{index.html} & HTML5 Semántico & Estructura del portal de inicio, bloque conceptual ODS 2 y 12, formulario de inserción y tabla de inventario. \\
\addlinespace
\texttt{js/nevera.js} & JavaScript (DOM) & Lógica de gestión de alimentos, alternancia de consejos (\texttt{style.display}), modo alerta de caducidad y reordenación (\texttt{insertBefore}). \\
\bottomrule
\end{tabularx}
\end{table}

\subsection*{3. Contribución al Trabajo Global (25\%)}
\begin{itemize}
    \item \textbf{Diseño de arquitectura general:} Creación del esqueleto principal del sitio web y homogeneización de la cabecera, navegación y pie común.
    \item \textbf{Gestión de alimentos:} Implementación del formulario de alta con validación de fechas y cálculo automático de la urgencia del consumo.
    \item \textbf{Reordenación y manipulación DOM:} Desarrollo de la funcionalidad de reordenación de productos urgentes en la cabecera de la tabla y eliminación dinámica con \texttt{removeChild}.
    \item \textbf{Dedicación estimada:} [Insertar número, ej. 25] horas de trabajo presencial y no presencial.
\end{itemize}

\subsection*{4. Declaración de Autoría y Compromiso Ético}
El estudiante abajo firmante declara que el código de los ficheros indicados ha sido desarrollado personalmente con la colaboración de su equipo, respetando la normativa de integridad académica de la UPV.

\vspace{1.2cm}
\begin{flushright}
    Firma del estudiante:\ \rule{6cm}{0.4pt}
\end{flushright}

\newpage

% ==============================================================================
% PÁGINA 2: LORA
% ==============================================================================
\begin{center}
    {\Large \textbf{UNIVERSITAT POLITÈCNICA DE VALÈNCIA}}\\[0.2cm]
    {\large Programación Web y Sistemas Inmersivos --- Curso 2026--2027}\\[0.3cm]
    {\LARGE \textbf{\color{azulOscuro}Documento de Relación de Trabajo Individual}}\\[0.2cm]
    {\small Proyecto: \textit{Tu Nevera Sostenible (ODS 2 y ODS 12)}}
\end{center}

\vspace{0.3cm}
\hrule
\vspace{0.4cm}

\subsection*{1. Datos del Estudiante}
\begin{tabularx}{\textwidth}{l X l X}
    \textbf{Nombre:} & Lora [Apellidos] & \textbf{DNI/NIE:} & [Completar DNI] \\
    \textbf{Email UPV:} & [lora@alumni.upv.es] & \textbf{Grupo:} & [Completar grupo] \\
\end{tabularx}

\vspace{0.4cm}
\subsection*{2. Ficheros Bajo Responsabilidad Única}
En cumplimiento de la normativa docente, cada fichero de código cuenta con un responsable exclusivo.
\begin{table}[h!]
\centering
\small
\begin{tabularx}{\textwidth}{p{4cm} p{3.5cm} X}
\toprule
\textbf{Fichero de Código} & \textbf{Tecnología} & \textbf{Descripción de Responsabilidad} \\
\midrule
\texttt{recetas.html} & HTML5 Semántico & Maquetación del área de juego de la olla, zona de despensa arrastrable y panel de resultados de recetas. \\
\addlinespace
\texttt{js/recetas.js} & JavaScript (HTML5 API) & API nativa de Drag and Drop (\texttt{dragstart}, \texttt{dragover}, \texttt{drop}), persistencia con \texttt{localStorage} y detección de movimiento de ratón. \\
\bottomrule
\end{tabularx}
\end{table}

\subsection*{3. Contribución al Trabajo Global (25\%)}
\begin{itemize}
    \item \textbf{Módulo de Drag and Drop:} Programación del flujo de arrastrar y soltar utilizando el objeto \texttt{dataTransfer} nativo de HTML5.
    \item \textbf{Persistencia en cliente:} Implementación de guardado y recuperación en \texttt{localStorage} para que los ingredientes en la cazuela se mantengan al recargar.
    \item \textbf{Interacción con ratón:} Configuración de eventos \texttt{mouseover}, \texttt{mousemove} y \texttt{mouseout} para mostrar información nutricional y coordenadas del cursor en vivo.
    \item \textbf{Motor culinario:} Lógica condicional de recetas de aprovechamiento tradicionales valencianas (arroz meloso, ropa vieja, torrijas).
    \item \textbf{Dedicación estimada:} [Insertar número, ej. 25] horas de trabajo presencial y no presencial.
\end{itemize}

\subsection*{4. Declaración de Autoría y Compromiso Ético}
El estudiante abajo firmante declara que el código de los ficheros indicados ha sido desarrollado personalmente con la colaboración de su equipo, respetando la normativa de integridad académica de la UPV.

\vspace{1.2cm}
\begin{flushright}
    Firma del estudiante:\ \rule{6cm}{0.4pt}
\end{flushright}

\newpage

% ==============================================================================
% PÁGINA 3: EDU
% ==============================================================================
\begin{center}
    {\Large \textbf{UNIVERSITAT POLITÈCNICA DE VALÈNCIA}}\\[0.2cm]
    {\large Programación Web y Sistemas Inmersivos --- Curso 2026--2027}\\[0.3cm]
    {\LARGE \textbf{\color{azulOscuro}Documento de Relación de Trabajo Individual}}\\[0.2cm]
    {\small Proyecto: \textit{Tu Nevera Sostenible (ODS 2 y ODS 12)}}
\end{center}

\vspace{0.3cm}
\hrule
\vspace{0.4cm}

\subsection*{1. Datos del Estudiante}
\begin{tabularx}{\textwidth}{l X l X}
    \textbf{Nombre:} & Edu [Apellidos] & \textbf{DNI/NIE:} & [Completar DNI] \\
    \textbf{Email UPV:} & [edu@alumni.upv.es] & \textbf{Grupo:} & [Completar grupo] \\
\end{tabularx}

\vspace{0.4cm}
\subsection*{2. Ficheros Bajo Responsabilidad Única}
En cumplimiento de la normativa docente, cada fichero de código cuenta con un responsable exclusivo.
\begin{table}[h!]
\centering
\small
\begin{tabularx}{\textwidth}{p{4cm} p{3.5cm} X}
\toprule
\textbf{Fichero de Código} & \textbf{Tecnología} & \textbf{Descripción de Responsabilidad} \\
\midrule
\texttt{impacto.html} & HTML5 Semántico & Estructura de la calculadora interactiva, controles deslizantes de entrada y tarjetas de métricas anuales. \\
\addlinespace
\texttt{js/impacto.js} & JavaScript (Cálculo) & Lógica matemática de conversión de alimentos a CO$_2$ evitado y ahorro en euros, y cambio dinámico de insignias ecológicas. \\
\bottomrule
\end{tabularx}
\end{table}

\subsection*{3. Contribución al Trabajo Global (25\%)}
\begin{itemize}
    \item \textbf{Modelado de datos ecológicos:} Investigación de factores de emisión de gases de efecto invernadero según categorías de alimentos (huerta, dieta mediterránea, carnes).
    \item \textbf{Procesamiento dinámico en cliente:} Actualización reactiva de valores en pantalla mediante eventos \texttt{input} en el slider de kilogramos.
    \item \textbf{Gamificación y estilos dinámicos:} Alteración de estilos, colores de fondo y textos del distintivo de nivel (Iniciado, Guardián, Campeón) según los kilos salvados.
    \item \textbf{Dedicación estimada:} [Insertar número, ej. 25] horas de trabajo presencial y no presencial.
\end{itemize}

\subsection*{4. Declaración de Autoría y Compromiso Ético}
El estudiante abajo firmante declara que el código de los ficheros indicados ha sido desarrollado personalmente con la colaboración de su equipo, respetando la normativa de integridad académica de la UPV.

\vspace{1.2cm}
\begin{flushright}
    Firma del estudiante:\ \rule{6cm}{0.4pt}
\end{flushright}

\newpage

% ==============================================================================
% PÁGINA 4: PABLO
% ==============================================================================
\begin{center}
    {\Large \textbf{UNIVERSITAT POLITÈCNICA DE VALÈNCIA}}\\[0.2cm]
    {\large Programación Web y Sistemas Inmersivos --- Curso 2026--2027}\\[0.3cm]
    {\LARGE \textbf{\color{azulOscuro}Documento de Relación de Trabajo Individual}}\\[0.2cm]
    {\small Proyecto: \textit{Tu Nevera Sostenible (ODS 2 y ODS 12)}}
\end{center}

\vspace{0.3cm}
\hrule
\vspace{0.4cm}

\subsection*{1. Datos del Estudiante}
\begin{tabularx}{\textwidth}{l X l X}
    \textbf{Nombre:} & Pablo [Apellidos] & \textbf{DNI/NIE:} & [Completar DNI] \\
    \textbf{Email UPV:} & [pablo@alumni.upv.es] & \textbf{Grupo:} & [Completar grupo] \\
\end{tabularx}

\vspace{0.4cm}
\subsection*{2. Ficheros Bajo Responsabilidad Única}
En cumplimiento de la normativa docente, cada fichero de código cuenta con un responsable exclusivo.
\begin{table}[h!]
\centering
\small
\begin{tabularx}{\textwidth}{p{4cm} p{3.5cm} X}
\toprule
\textbf{Fichero de Código} & \textbf{Tecnología} & \textbf{Descripción de Responsabilidad} \\
\midrule
\texttt{donaciones.html} & HTML5 Semántico & Directorio de comedores sociales de València, formulario de registro de donaciones y marco en línea \texttt{<iframe>}. \\
\addlinespace
\texttt{mapa-valencia.html} & HTML5 / Leaflet & Implementación del mapa interactivo con OpenStreetMap para mostrar los 4 puntos simultáneos y enfoque individual. \\
\addlinespace
\texttt{js/donaciones.js} & JavaScript (Eventos) & Validación de campos de formulario, confirmación dinámica en pantalla y filtrado de tablas por distritos. \\
\bottomrule
\end{tabularx}
\end{table}

\subsection*{3. Contribución al Trabajo Global (25\%)}
\begin{itemize}
    \item \textbf{Integración del mapa interactivo:} Configuración del documento para el \texttt{<iframe>} asegurando visualización de los 4 centros simultáneos y zoom interactivo al pulsar.
    \item \textbf{Filtrado dinámico de tablas:} Programación del evento \texttt{change} sobre el selector de zonas para alternar la visibilidad (\texttt{display: none}) de las filas de centros.
    \item \textbf{Gestión de donaciones solidarias:} Desarrollo del formulario con \texttt{<fieldset>} y confirmación en cliente tras aceptar la condición higiénico-sanitaria.
    \item \textbf{Dedicación estimada:} [Insertar número, ej. 25] horas de trabajo presencial y no presencial.
\end{itemize}

\subsection*{4. Declaración de Autoría y Compromiso Ético}
El estudiante abajo firmante declara que el código de los ficheros indicados ha sido desarrollado personalmente con la colaboración de su equipo, respetando la normativa de integridad académica de la UPV.

\vspace{1.2cm}
\begin{flushright}
    Firma del estudiante:\ \rule{6cm}{0.4pt}
\end{flushright}

\end{document}
